\documentclass[11pt]{article}
\usepackage[a4paper,margin=28mm]{geometry}
\usepackage{amsmath,amssymb,amsthm,hyperref}
\newtheorem{proposition}{Proposition}
\title{Many polynomial-size dice would settle the exponential-total question}
\author{Asymptotic reductions from hereditary fairness}
\date{30 September 2026}
\begin{document}\maketitle
Let $A_r$ be the minimum total number of faces of a permutation-fair
family of $r$ finite equiprobable dice, with arbitrary individual sizes.
All logarithms below are natural.

\begin{proposition}[Arithmetic upper bound on the exceptional set]
If a permutation-fair family contains $r$ dice with at most $M$ faces
each, then, as $r\to\infty$,
\[
 \log M\ge(1/2+o(1))r.
\]
Consequently an $N$-die family contains at most
$(2C+o(1))\log N$ dice with at most $N^C$ faces, for fixed $C>0$.
\end{proposition}
\begin{proof}
Restrict to the $r$ specified dice, with counts $K_1,\ldots,K_r$.
For every prime $p\le r$, at most $p-1$ of these counts can fail to
be divisible by $p$: otherwise a $p$-element restriction would have
a face-count product not divisible by $p!$. Thus
\[
 r\log M\ge\sum_{i=1}^r\log K_i
 \ge\sum_{p\le r}(r-p+1)\log p.
\]
Writing $\vartheta(x)=\sum_{p\le x}\log p$, the last sum is
$\vartheta(r)+\int_2^r\vartheta(t)\,dt=(1/2+o(1))r^2$
by the prime number theorem. This proves the claims; a bounded $r$
is harmless in the second asymptotic assertion.
\end{proof}

\begin{proposition}[Reduction to the first open problem]\label{reduction}
Suppose there are absolute constants $c,C>0$ such that, for every
sufficiently large integer $N$, some permutation-fair $N$-die family
contains at least $c\log N$ dice with at most $N^C$ faces each.
Then
\[
 A_r\le\exp(O(r))\qquad(r\to\infty).
\]
More precisely $\log A_r\le(C/c+o(1))r$.
\end{proposition}
\begin{proof}
For each large $r$, take $N=\lceil\exp((r+1)/c)\rceil$ and retain
any $r$ of the promised small dice. Hereditary fairness gives a fair
$r$-die family of total size at most $rN^C$. Its logarithm is
at most $(C/c)r+O(\log r)$.
\end{proof}

\paragraph{Interpretation and limits.}
The first proposition is a sharpening of the manuscript's arithmetic
maximum-die argument, not a new construction. The second is a one-way
reduction. A proof of logarithmically many uniformly polynomial-size
exceptional dice would already close the original exponential versus
$\exp(O(r\log r))$ upper-bound gap. Conversely, an exponential-total
construction for $r$ dice does not automatically embed into a much larger
fair family while keeping those dice small; no equivalence is claimed.

The proved fixed-$k$ construction, including its uniform quantitative
bound
\[
 S\le(N+2^k+2)^{\exp(Ck^2)}
\]
for $k$ specified equal-count dice, does not meet the uniform hypotheses
of Proposition~\ref{reduction}. Its growing-$k$ consequences give
subexponential sizes, not a bound $N^C$ with an absolute exponent.
See \path{size_bounds/uniform_fixed_number_complexity.tex}.
\end{document}
